\documentclass{fiche}
\metadonnees{12}{Sous la salle à manger}{Bloc 3 — La ville, le travail, la science}

\begin{document}
\entete

\objectifs{%
\textbf{Langue} : dénombrables et indénombrables ; quantifieurs ; \emph{some} et
\emph{any}.\par
\textbf{Texte} : George Orwell, \emph{Down and Out in Paris and London} (1933), chapitres
\textsc{x} et \textsc{xii}.\par
\textbf{Durée} : 1\,h\,30 — exercices 35 min, texte et questions 30 min, version 25 min.}

%% ===================================================================
\exercice[13 min]{Dénombrable ou indénombrable}

\rappel{\textbf{Rappel.} Un indénombrable n'a pas de pluriel et refuse \emph{a} / \emph{an} :
pour le compter, on passe par une unité (\emph{a piece of}, \emph{a lump of}). Il est suivi
du verbe au singulier.}

\consigne{Complétez avec \emph{a}, \emph{an}, \emph{some} ou rien (Ø), et mettez le nom au
pluriel s'il le faut.}

\begin{anglais}
\begin{enumerate}[label=\arabic*.,itemsep=2.2mm,leftmargin=*]
  \item We had to fetch \trou{some} bread and \trou{some} coffee from the kitchen.
  \item He counted \trou{Ø} lumps of sugar all morning.
  \item There was \trou{a} thermometer in the cellar.
  \item The room had \trou{a} dirty smell of food and sweat.
  \item \trou{Ø} work in the cafeterie came in bursts of two hours.
  \item They stole \trou{Ø} food and hid it behind the crockery.
  \item He gave me \trou{a} piece of advice I have never forgotten.
  \item \trou{Ø} furniture in the cellar was broken.
  \item We had no time to sweep \trou{the} floor before evening.
  \item There was \trou{Ø} scarcely any news from the dining-room.
  \item She had \trou{Ø} beautiful hair, but it was \trou{a} wig.
  \item He asked me for \trou{some} information about the hotel.
  \item The waiters brought back \trou{Ø} dirty crockery by the armful.
  \item I could not bear \trou{the} smoke of the gas-fires.
\end{enumerate}
\end{anglais}

\note[Huit indénombrables qui trompent les francophones]{\emph{advice} (un conseil = \emph{a
piece of advice}), \emph{information}, \emph{news} (singulier malgré le \emph{-s} :
\emph{the news is good}), \emph{furniture}, \emph{luggage}, \emph{work}, \emph{progress},
\emph{hair} (la chevelure ; \emph{a hair} = un cheveu isolé). \emph{Crockery} et
\emph{machinery} appartiennent à la même famille : des collectifs en \emph{-ery}.}

%% ===================================================================
\exercice[10 min]{Quantifieurs}

\rappel{\textbf{Rappel.} \emph{Few} et \emph{little} disent le manque ; \emph{a few} et
\emph{a little} disent qu'il y en a un peu. \emph{Many} et \emph{few} comptent,
\emph{much} et \emph{little} mesurent.}

\consigne{Complétez avec \emph{much}, \emph{many}, \emph{a lot of}, \emph{few}, \emph{a few},
\emph{little} ou \emph{a little}.}

\begin{anglais}
\begin{enumerate}[label=\arabic*.,itemsep=2.2mm,leftmargin=*]
  \item There was \trou{little} time to clean the floor. \emph{(il n'y en avait presque pas)}
  \item Give me \trou{a little} more water, please.
  \item \trou{Few} customers ever saw the kitchens. \emph{(presque aucun)}
  \item \trou{A few} waiters were already at the table, mixing salads.
  \item How \trou{much} work do you do in a week?
  \item How \trou{many} hours a day did she stand at the sink?
  \item There were \trou{a lot of} plates and only two sinks.
  \item He spoke \trou{little} English and no French at all.
  \item We had \trou{a lot of} trouble getting common civility from them.
  \item \trou{Many} of the men had been there for years.
\end{enumerate}
\end{anglais}

\note[Trois remarques]{\emph{Much} et \emph{many} appartiennent surtout aux phrases
négatives et interrogatives ; à l'affirmative on préfère \emph{a lot of} ou \emph{plenty
of} (7, 9). Après \emph{how}, en revanche, \emph{much} et \emph{many} sont obligatoires (5,
6). Enfin \emph{few} et \emph{little} sans article ont une valeur négative : \emph{few
customers} signifie qu'il n'en venait presque aucun, \emph{a few customers} qu'il en venait
quelques-uns.}

%% ===================================================================
\exercice[12 min]{La cuisine et le service}
\consigne{a) Associez chaque mot à sa traduction.}

\begin{multicols}{2}
\begin{enumerate}[label=\arabic*.,itemsep=1.2mm,leftmargin=*]
  \item a scullery \hfill \trou{f}
  \item crockery \hfill \trou{b}
  \item a sink \hfill \trou{i}
  \item a cellar \hfill \trou{a}
  \item a cupboard \hfill \trou{k}
  \item a lift \hfill \trou{d}
  \item to rinse \hfill \trou{h}
  \item to sweep \hfill \trou{c}
  \item a tray \hfill \trou{l}
  \item a shift \hfill \trou{e}
  \item filth \hfill \trou{g}
  \item sweat \hfill \trou{j}
\end{enumerate}
\columnbreak
\begin{enumerate}[label=\alph*.,itemsep=1.2mm,leftmargin=*]
  \item une cave
  \item la vaisselle
  \item balayer
  \item un monte-charge, un ascenseur
  \item une équipe, un poste (de travail)
  \item une plonge, une arrière-cuisine
  \item la crasse
  \item rincer
  \item un évier
  \item la sueur
  \item un placard
  \item un plateau
\end{enumerate}
\end{multicols}

\note[Deux collectifs]{\emph{Crockery}, \emph{cutlery}, \emph{machinery}, \emph{jewellery} :
le suffixe \emph{-ery} forme des noms collectifs indénombrables. On dit \emph{a piece of
crockery}, jamais \emph{a crockery}. \emph{Filth} et \emph{sweat} sont indénombrables
également : la crasse et la sueur ne se comptent pas.}

\consigne{b) Complétez avec un mot de la liste.}
\begin{anglais}
\begin{enumerate}[label=\arabic*.,itemsep=2mm,leftmargin=*]
  \item There were only two \trou{sinks} and no washing basin.
  \item Nobody had time to \trou{sweep} the floor before the evening.
  \item Behind the piles of \trou{crockery} the waiters hid stolen food.
  \item The \trou{cellar} was so hot that the temperature scarcely fell below 110 degrees.
\end{enumerate}
\end{anglais}

%% ===================================================================
\newpage
\section{Texte}

\noindent\small\textit{À vingt-cinq ans, Orwell vit à Paris sans argent et travaille comme
plongeur dans les sous-sols d'un grand hôtel qu'il appelle l'Hôtel X.}\normalsize

\begin{texte}
Our \gl[a cafeterie]{cafeterie}{l'office (mot français gardé tel quel)} was a murky cellar
measuring twenty feet by seven by eight high, and so crowded with coffee-urns, breadcutters
and the like that one could hardly move without banging against something. It was lighted by
one dim electric bulb, and four or five gas-fires that sent out a fierce red breath. There
was a thermometer there, and the temperature never fell below 110 degrees Fahrenheit --- it
neared 130 at some times of the day. At one end were five service lifts, and at the other an
ice cupboard where we stored milk and butter. When you went into the ice cupboard you dropped
a hundred degrees of temperature at a single step; it used to remind me of the hymn about
Greenland's icy mountains and India's coral strand. Two men worked in the cafeterie besides
Boris and myself. One was Mario, a huge, excitable Italian --- he was like a city policeman
with operatic gestures --- and the other, a hairy, \gl[uncouth]{uncouth}{fruste, mal
dégrossi} animal whom we called the Magyar; I think he was a Transylvanian, or something
even more remote. Except the Magyar we were all big men, and at the rush hours we collided
incessantly.

[\ldots]

It was amusing to look round the filthy little \gl[a scullery]{scullery}{la plonge,
l'arrière-cuisine} and think that only a double door was between us and the dining-room.
There sat the customers in all their splendour --- spotless table-cloths, bowls of flowers,
mirrors and \gl[a gilt cornice]{gilt cornices}{une corniche dorée} and painted
\gl[cherubim]{cherubim}{des chérubins} ; and here, just a few feet away, we in our disgusting
\gl[filth]{filth}{la crasse}. For it really was disgusting filth. There was no time to sweep
the floor till evening, and we \gl[to slither]{slithered}{glisser, patiner} about in a
compound of soapy water, lettuce-leaves, torn paper and trampled food. A dozen waiters with
their coats off, showing their sweaty \gl[an armpit]{armpits}{une aisselle}, sat at the table
mixing salads and sticking their thumbs into the cream pots. The room had a dirty, mixed
smell of food and sweat. Everywhere in the cupboards, behind the piles of
\gl[crockery]{crockery}{la vaisselle}, were \gl[squalid]{squalid}{sordide} stores of food
that the waiters had stolen. There were only two \gl[a sink]{sinks}{un évier}, and no washing
basin, and it was nothing unusual for a waiter to wash his face in the water in which clean
crockery was rinsing. But the customers saw nothing of this. There were a coco-nut mat and a
mirror outside the dining-room door, and the waiters used to \gl[to preen oneself]{preen
themselves}{se pomponner, se lisser les plumes} up and go in looking the picture of
cleanliness.

It is an instructive sight to see a waiter going into a hotel dining-room. As he passes the
door a sudden change comes over him. The set of his shoulders alters; all the dirt and hurry
and irritation have dropped off in an instant. He \gl[to glide]{glides}{glisser} over the
carpet, with a solemn priest-like air. [\ldots] Then he entered the dining-room and sailed
across it dish in hand, graceful as a swan. Ten seconds later he was bowing
\gl[reverently]{reverently}{avec déférence} to a customer. And you could not help thinking, as
you saw him bow and smile, with that \gl[benign]{benign}{bienveillant} smile of the trained
waiter, that the customer was put to shame by having such an aristocrat to serve him.

This washing up was a thoroughly \gl[odious]{odious}{odieux, détestable} job --- not hard, but
boring and silly beyond words. It is dreadful to think that some people spend their whole
decades at such occupations.
\end{texte}
\source{George Orwell, \emph{Down and Out in Paris and London}, Londres, Victor Gollancz,
1933, ch.~\textsc{x} et \textsc{xii}.}

\partiel[15 min]{Questions}
\consigne{\en{Answer in English, in full sentences.}}

\begin{enumerate}[label=\arabic*.,itemsep=2mm,leftmargin=*]
  \item \en{Describe the cafeterie: size, light, heat.}
    \lignes{3}
    \corr{A murky cellar twenty feet by seven and eight high, so crowded that one could
    hardly move; lit by a single dim bulb and four or five gas-fires; never below 110 degrees
    Fahrenheit and sometimes near 130.}
  \item \en{What separates the kitchens from the dining-room, and what is on each side?}
    \lignes{3}
    \corr{Only a double door. On one side, spotless table-cloths, flowers, mirrors, gilt
    cornices and painted cherubim; on the other, a few feet away, soapy water, lettuce
    leaves, torn paper, trampled food and men in disgusting filth.}
  \item \en{What happens to a waiter as he passes through the door?}
    \lignes{3}
    \corr{A sudden change comes over him: his shoulders alter, all the dirt, hurry and
    irritation drop off, and he glides over the carpet with a solemn, priest-like air.}
  \item \en{Why does Orwell compare the waiter to a priest, a swan and an aristocrat?}
    \lignes{3}
    \corr{All three suggest a dignity that is entirely performed. The comparisons raise him
    as high as possible precisely because we have just seen where he comes from, and the gap
    is the point.}
  \item \en{``The customer was put to shame by having such an aristocrat to serve him.''
    What is ironic here?}
    \lignes{3}
    \corr{The man who is humiliated in reality is the waiter; in appearance it is the
    customer who is outclassed. The social order is exactly reversed in the twenty seconds of
    the performance, and only in those.}
  \item \en{What does Orwell say about the work itself?}
    \lignes{2}
    \corr{That it was not hard but boring and silly beyond words, and that it is dreadful to
    think of people spending whole decades at such occupations.}
\end{enumerate}

%% ===================================================================
\newpage
\section{Version}
\consigne{Traduisez le passage en français. Il suit la description de l'office.}

\begin{version}
The work in the cafeterie was \gl[spasmodic]{spasmodic}{par à-coups}. We were never idle, but
the real work only came in bursts of two hours at a time --- we called each burst
\emph{un coup de feu}. The first coup de feu came at eight, when the guests upstairs began to
wake up and demand breakfast. At eight a sudden banging and yelling would break out all
through the basement; bells rang on all sides, blue-\gl[aproned]{aproned}{en tablier} men
rushed through the passages, our service lifts came down with a simultaneous crash, and the
waiters on all five floors began shouting Italian \gl[an oath]{oaths}{un juron} down the
\gl[a shaft]{shafts}{une cage, un conduit}. I don't remember all our duties, but they
included making tea, coffee and chocolate, fetching meals from the kitchen, wines from the
cellar and fruit and so forth from the dining-room, slicing bread, making toast, rolling
\gl[a pat of butter]{pats of butter}{une plaquette de beurre}, measuring jam, opening
milk-cans, counting \gl[a lump of sugar]{lumps of sugar}{un morceau de sucre}, boiling eggs,
cooking \gl[porridge]{porridge}{la bouillie d'avoine}, pounding ice, grinding coffee --- all
this for from a hundred to two hundred customers.
\end{version}
\source{\emph{Ibid.}}

\lignes{22}

\corr{\textbf{Proposition de traduction.} Le travail à l'office se faisait par à-coups. Nous
n'étions jamais oisifs, mais la vraie besogne ne venait que par accès de deux heures d'affilée
--- nous appelions chacun de ces accès un coup de feu. Le premier coup de feu arrivait à huit
heures, quand les clients, en haut, commençaient à se réveiller et à réclamer leur petit
déjeuner. À huit heures, un vacarme soudain de coups et de cris éclatait dans tout le
sous-sol ; les sonnettes retentissaient de tous côtés, des hommes en tablier bleu se
précipitaient dans les couloirs, nos monte-charges descendaient tous ensemble dans un
fracas, et les garçons des cinq étages se mettaient à hurler des jurons italiens dans les
cages. Je ne me rappelle pas toutes nos tâches, mais elles comprenaient : faire le thé, le
café et le chocolat, rapporter les plats de la cuisine, les vins de la cave, les fruits et
le reste de la salle à manger, couper le pain en tranches, faire des toasts, rouler des
plaquettes de beurre, doser la confiture, ouvrir les bidons de lait, compter les morceaux de
sucre, faire cuire les œufs, préparer la bouillie d'avoine, piler la glace, moudre le café
--- tout cela pour cent à deux cents clients.}

\note[Choix de traduction 1]{\emph{a sudden banging and yelling would break out} : ce
\emph{would} n'est pas un conditionnel mais l'habitude passée de la fiche 2. Tout le
paragraphe décrit ce qui se passait chaque matin : imparfait en français.}

\note[Choix de traduction 2]{\emph{making tea, coffee and chocolate, fetching meals\ldots} :
la longue liste de gérondifs se rend par des infinitifs, plus naturels en français après
\og comprendre \fg. Garder la liste intacte : son accumulation est le sens même du
passage.}

\note[Choix de traduction 3]{\emph{slicing bread}, \emph{rolling pats of butter},
\emph{pounding ice}, \emph{grinding coffee} : à chaque fois un verbe anglais unique et précis
là où le français demande un verbe plus un complément (\og couper en tranches \fg). C'est le
contraire de l'étoffement des adjectifs : ici c'est le verbe qui doit se déplier.}

\note[Choix de traduction 4]{\emph{all this for from a hundred to two hundred customers} :
deux prépositions se suivent, ce que le français interdit. On supprime la première : \og tout
cela pour cent à deux cents clients \fg.}

%% ===================================================================
\partiel{Lexique à mémoriser}
\begin{itemize}[label=--,itemsep=0.8mm,leftmargin=*]\small
  \item \textbf{a cellar} — une cave ; \textit{et} a basement \textit{« un sous-sol »}
  \item \textbf{murky} — sombre, glauque
  \item \textbf{dim} — faible, blafard ; a dim light
  \item \textbf{crowded} — bondé ; \textit{nom} a crowd
  \item \textbf{a sink} — un évier
  \item \textbf{crockery} — la vaisselle \textit{(indénombrable en} -ery\textit{)}
  \item \textbf{cutlery} — les couverts \textit{(même formation)}
  \item \textbf{filth} — la crasse ; \textit{adjectif} filthy
  \item \textbf{sweat} — la sueur ; \textit{adjectif} sweaty\textit{, se prononce} swet
  \item \textbf{to sweep} — balayer ; swept, swept
  \item \textbf{to rinse} — rincer
  \item \textbf{a duty} — une tâche, un devoir ; \textit{adjectif} dutiful
  \item \textbf{a burst} — un accès, une explosion ; \textit{verbe} to burst \textit{« éclater »}
  \item \textbf{idle} — oisif ; \textit{nom} idleness
  \item \textbf{to demand} — exiger ; \textit{et non « demander »} --- \textit{faux ami :} to ask
  \item \textbf{to fetch} — aller chercher
  \item \textbf{to slice} — couper en tranches ; \textit{nom} a slice
  \item \textbf{to grind} — moudre ; ground, ground
  \item \textbf{to bow} — s'incliner ; \textit{rime avec} now
  \item \textbf{to put to shame} — faire honte à ; \textit{adjectif} ashamed
\end{itemize}

\end{document}
